\documentclass[../../main.tex]{subfiles}

\begin{document}

\section{紧度量空间}\label{section:ch09-紧度量空间}

回忆集族 $\{E_{\lambda}\}_{\lambda\in\Lambda}$ 称为一个集合 $E$ 的覆盖, 若 $E\subseteq\bigcup_{\lambda\in\Lambda}E_{\lambda}$. 谈到 $E$ 的覆盖的子覆盖我们指的是该覆盖的子族自身也是 $E$ 的覆盖. 若 $E$ 是度量空间 $X$ 的子集, 谈到 $E$ 的开覆盖我们指由 $X$ 的开子集组成的 $E$ 的覆盖. 第 1 章讨论过的关于实数集的紧性的概念推广到如下度量空间类.

\begin{definition}[紧]\label{definition:ch09-紧}
度量空间 $X$ 称为紧的, 若 $X$ 的每个开覆盖有有限子覆盖. $X$ 的子集 $K$ 称为紧的, 若视为 $X$ 的度量子空间的 $K$ 是紧的.
\end{definition}

度量空间 $X$ 的子空间 $K$ 的开子集是 $K$ 与 $X$ 的开集的交. 因此一个度量空间 $X$ 的子集 $K$ 是紧的, 当且仅当 $K$ 的每个由 $X$ 的开子集族组成的覆盖有有限的子覆盖.

若 $\mathcal{T}$ 是度量空间 $X$ 的开子集族, 则 $\mathcal{T}$ 中的集合的补集族 $\mathcal{F}$ 是闭集的族. 此外, $\mathcal{T}$ 是一个覆盖当且仅当 $\mathcal{F}$ 有空的交. 因此, 根据 De Morgan 等式, 一个度量空间 $X$ 是紧的当且仅当每个具有空交集的闭集族有一个有限子族, 其交也是空的. $X$ 中的集族 $\mathcal{F}$ 称为具有有限交性质, 若 $\mathcal{F}$ 的任何有限子族有一个非空的交集. 因此我们可以用闭集族叙述紧性如下.

\begin{proposition}\label{proposition:ch09-紧有限交}
度量空间 $X$ 是紧的, 当且仅当 $X$ 的每个具有有限交性质的闭集族 $\mathcal{F}$ 有非空的交集.
\end{proposition}

\begin{definition}[全有界]\label{definition:ch09-全有界}
度量空间 $X$ 称为全有界的, 若对每个 $\varepsilon>0$, 空间 $X$ 可被有限个半径为 $\varepsilon$ 的开球覆盖. $X$ 的子集 $E$ 称为全有界, 若考虑 $E$ 为度量空间 $X$ 的子空间, $E$ 是全有界的.
\end{definition}

对于度量空间 $X$ 的子集 $E$, 谈到 $E$ 的 $\varepsilon$-网我们指的是有限个开球族 $\{B(x_{k},\varepsilon)\}_{k=1}^{n}$, 其中心 $x_{k}$ 属于 $X$ 并覆盖 $E$. 我们把证明每个度量空间 $E$ 是全有界的, 当且仅当对每个 $\varepsilon>0$ 存在 $E$ 的 $\varepsilon$-网留作练习. 该观察的要点在于度量子空间 $E$ 为全有界的准则不必要求该网中的球的中心属于 $E$.

若一个度量空间 $X$ 是全有界的, 则它在直径有限的意义下是有界的. 的确, 若 $X$ 被有限个半径为 $1$ 的开球覆盖, 则从三角不等式推出 $\mathrm{diam}X\leqslant c$, 其中 $c=2+d$, $d$ 是这些覆盖球中心之间的最大距离. 然而, 从下面的例子看到, 一个有界度量空间不必是全有界的.

\begin{example}\label{example:ch09-l2单位球不全有界}
令 $X$ 为平方可和序列的 Banach 空间 $\ell^{2}$. 考虑闭单位球 $B=\{\{x_{n}\}\in\ell^{2}\mid\|\{x_{n}\}\|_{2}\leqslant 1\}$. 则 $B$ 是有界的. 我们宣称 $B$ 不是全有界的. 事实上, 对每个自然数 $n$, 令 $e_{n}$ 为第 $n$ 个分量为 $1$ 而其余分量为 $0$ 的序列. 若 $m\neq n$, 则 $\|e_{n}-e_{m}\|_{2}=\sqrt{2}$. 因而 $B$ 不可能包含于有限个半径为 $r<1/2$ 的球, 这是由于这些球中的一个含 $\{e_{n}\}$ 的两个元素, 它们相距 $\sqrt{2}$ 但球的直径小于 $1$.
\end{example}

\begin{proposition}\label{proposition:ch09-Rn有界全有界}
欧氏空间 $\mathbf{R}^{n}$ 的子集是有界的当且仅当它是全有界的.
\end{proposition}

\begin{proof}
全有界度量空间总是有界的. 因此, 令 $E$ 为 $\mathbf{R}^{n}$ 的有界子集. 为简单起见取 $n=2$. 令 $\varepsilon>0$. 由于 $E$ 是有界的, 我们可以取 $a>0$ 足够大以使得 $E$ 包含于正方形 $[-a,a]\times[-a,a]$ 中. 令 $P_{k}$ 为 $[-a,a]$ 的分划, 使得每个分割后的区间的长度小于 $1/k$. 则 $P_{k}\times P_{k}$ 诱导了 $[-a,a]\times[-a,a]$ 的一个直径至多为 $\sqrt{2}/k$ 的闭矩形的分划. 选取 $k$ 使得 $\sqrt{2}/k<\varepsilon$. 考虑中心在 $(x,y)$、半径为 $\varepsilon$ 的有限球族, 其中 $x$ 和 $y$ 是 $P_{k}$ 的分点. 则该半径为 $\varepsilon$ 的有限球族覆盖了正方形 $[-a,a]\times[-a,a]$, 因此也覆盖了 $E$.
\end{proof}

\begin{definition}[序列紧]\label{definition:ch09-序列紧}
度量空间 $X$ 称为序列紧的, 若 $X$ 中的每个序列有收敛到 $X$ 中的点的子序列.
\end{definition}

\begin{theorem}[紧性的刻画]\label{theorem:ch09-紧性刻画}
对于度量空间 $X$, 以下三个断言等价:
\begin{enumerate}[(i)]
\item $X$ 是完备的与全有界的;
\item $X$ 是紧的;
\item $X$ 是序列紧的.
\end{enumerate}
\end{theorem}

为清晰起见, 我们将证明分为三个命题.

\begin{proposition}\label{proposition:ch09-完备全有界紧}
若度量空间 $X$ 是完备的与全有界的, 则它是紧的.
\end{proposition}

\begin{proof}
我们用反证法讨论. 假定 $\{\mathcal{O}_{\lambda}\}_{\lambda\in\Lambda}$ 是 $X$ 的开覆盖, 其没有有限子覆盖. 由于 $X$ 是全有界的, 我们可以选取覆盖 $X$ 的有限个半径小于 $1/2$ 的开球族. 这些球必有一个不能被 $\{\mathcal{O}_{\lambda}\}_{\lambda\in\Lambda}$ 的有限子族覆盖. 选取这样一个球且将它的闭包记为 $F_{1}$. 则 $F_{1}$ 是闭的且 $\mathrm{diam}F_{1}\leqslant 1$. 又一次用 $X$ 的全有界性, 存在覆盖 $X$ 的有限个半径小于 $1/4$ 的开球族. 该族也覆盖 $F_{1}$. 这些球必有一个不能被 $\{\mathcal{O}_{\lambda}\}_{\lambda\in\Lambda}$ 的有限子族覆盖. 定义 $F_{2}$ 为这样的一个球与 $F_{1}$ 的交的闭包. 则 $F_{1}$ 和 $F_{2}$ 是闭的, $F_{2}\subseteq F_{1}$, 且 $\mathrm{diam}F_{1}\leqslant 1$, $\mathrm{diam}F_{2}\leqslant 1/2$. 以这种方式继续, 我们得到非空闭子集的收缩序列 $\{F_{n}\}$, 它具有每个 $F_{n}$ 不能被 $\{\mathcal{O}_{\lambda}\}_{\lambda\in\Lambda}$ 的有限子族覆盖的性质. 但 $X$ 是完备的. 根据 Cantor 交定理, 存在 $X$ 中的点 $x_{0}$ 属于交 $\bigcap_{n=1}^{\infty}F_{n}$. 存在某个 $\lambda_{0}$ 使得 $\mathcal{O}_{\lambda_{0}}$ 包含 $x_{0}$, 且因为 $\mathcal{O}_{\lambda_{0}}$ 是开的, 存在一个中心在 $x_{0}$ 的球 $B(x_{0},r)$, 使得 $B(x_{0},r)\subseteq\mathcal{O}_{\lambda_{0}}$. 因为 $\lim_{n\to\infty}\mathrm{diam}F_{n}=0$ 且 $x_{0}\in\bigcap_{n=1}^{\infty}F_{n}$, 则有一个指标 $n$ 使得 $F_{n}\subseteq\mathcal{O}_{\lambda_{0}}$. 这与 $F_{n}$ 不能被 $\{\mathcal{O}_{\lambda}\}_{\lambda\in\Lambda}$ 的有限子族覆盖矛盾. 这个矛盾表明 $X$ 是紧的.
\end{proof}

\begin{proposition}\label{proposition:ch09-紧序列紧}
若度量空间 $X$ 是紧的, 则它是序列紧的.
\end{proposition}

\begin{proof}
令 $\{x_{n}\}$ 为 $X$ 中的序列. 对每个指标 $n$, 令 $F_{n}$ 为非空集 $\{x_{k}\mid k\geqslant n\}$ 的闭包. 则 $\{F_{n}\}$ 是递减的非空闭集序列. 根据 Cantor 交定理, 存在 $X$ 中的 $x_{0}$ 属于交 $\bigcap_{n=1}^{\infty}F_{n}$. 由于对每个 $n$, $x_{0}$ 属于 $\{x_{k}\mid k\geqslant n\}$ 的闭包, 球 $B(x_{0},1/k)$ 与 $\{x_{k}\mid k\geqslant n\}$ 有非空交. 根据归纳我们可以选取一个严格递增的指标序列 $\{n_{k}\}$, 使得对每个指标 $k$ 有 $\rho(x_{0},x_{n_{k}})<1/k$. 子序列 $\{x_{n_{k}}\}$ 收敛于 $x_{0}$. 因此 $X$ 是序列紧的.
\end{proof}

\begin{proposition}\label{proposition:ch09-序列紧完备全有界}
若度量空间 $X$ 是序列紧的, 则它是完备与全有界的.
\end{proposition}

\begin{proof}
我们用反证法证明全有界性. 假定 $X$ 不是全有界的. 则对某个 $\varepsilon>0$, 我们不能用有限个半径为 $\varepsilon$ 的开球覆盖 $X$. 选取 $X$ 中的点 $x_{1}$. 由于 $X$ 不包含于 $B(x_{1},\varepsilon)$, 我们可以选取 $x_{2}\in X$ 使得 $\rho(x_{1},x_{2})\geqslant\varepsilon$. 现在 $X$ 不包含于 $B(x_{1},\varepsilon)\bigcup B(x_{2},\varepsilon)$, 我们可以选取 $x_{3}\in X$ 使得 $\rho(x_{3},x_{2})\geqslant\varepsilon$, $\rho(x_{3},x_{1})\geqslant\varepsilon$. 以这种方式继续, 我们得到 $X$ 中具有当 $n>k$ 时 $\rho(x_{n},x_{k})\geqslant\varepsilon$ 的性质的序列 $\{x_{n}\}$. 则序列 $\{x_{n}\}$ 没有收敛的子序列, 这是由于任何子序列的任何两个不同的项距离 $\varepsilon$ 或更远. 因此 $X$ 不是序列紧的. 这个矛盾表明 $X$ 必须是全有界的.

为证明 $X$ 是完备的, 令 $\{x_{n}\}$ 为 $X$ 中的 Cauchy 序列. 由于 $X$ 是序列紧的, $\{x_{n}\}$ 的子序列收敛到点 $x\in X$. 利用 Cauchy 性质不难证明整个序列收敛到 $x$. 因此 $X$ 是完备的.
\end{proof}

这三个命题完成了紧性定理的证明.

由于欧氏空间 $\mathbf{R}^{n}$ 是完备的, 作为度量子空间每个闭子集是完备的. 此外, 命题 15 断言欧氏空间的子集是有界的当且仅当它是全有界的. 因此从紧性定理我们有以下欧氏空间的子空间的紧性刻画.

\begin{theorem}\label{theorem:ch09-Rn紧性刻画}
对于 $\mathbf{R}^{n}$ 的子集 $K$, 以下三个断言等价:
\begin{enumerate}[(i)]
\item $K$ 是闭且有界的.
\item $K$ 是紧的.
\item $K$ 是序列紧的.
\end{enumerate}
\end{theorem}

关于该定理, (i) 与 (ii) 的等价性即是已知的 Heine-Borel 定理, 而 (i) 与 (iii) 的等价性即是 Bolzano-Weierstrass 定理. 在第 1 章, 我们在 $\mathbf{R}=\mathbf{R}^{1}$ 情形证明了它们, 因为我们在建立实变量函数的 Lebesgue 积分时用到它们.

\begin{proposition}\label{proposition:ch09-紧集连续象紧}
令 $f$ 为从紧度量空间 $X$ 到度量空间 $Y$ 的连续映射. 则它的象 $f(X)$ 也是紧的.
\end{proposition}

\begin{proof}
令 $\{\mathcal{O}_{\lambda}\}_{\lambda\in\Lambda}$ 为 $f(X)$ 的一个开覆盖. 则根据 $f$ 的连续性, $\{f^{-1}(\mathcal{O}_{\lambda})\}_{\lambda\in\Lambda}$ 是 $X$ 的一个开覆盖. 根据 $X$ 的紧性, 也存在覆盖 $X$ 的有限子族 $\{f^{-1}(\mathcal{O}_{\lambda_{1}}),\cdots,f^{-1}(\mathcal{O}_{\lambda_{n}})\}$. 由于 $f$ 将 $X$ 映上 $f(X)$, 有限族 $\{\mathcal{O}_{\lambda_{1}},\cdots,\mathcal{O}_{\lambda_{n}}\}$ 覆盖了 $f(X)$.
\end{proof}

在微积分课程以及我们第 1 章证明过的一元实变函数的首要性质之一是, 闭有界区间上的连续函数取到最大值与最小值. 我们自然尝试对那些使得该极值性质成立的度量空间归类.

\begin{theorem}\label{theorem:ch09-紧集连续函数取极值}
令 $X$ 为度量空间. 则 $X$ 是紧的当且仅当 $X$ 上的每个连续实值函数取到最大值与最小值.
\end{theorem}

\begin{proof}
首先假定 $X$ 是紧的. 令函数 $f:X\to\mathbf{R}$ 为连续的. 前一个命题告诉我们 $f(X)$ 是实数的紧集. 根据定理 20, $f(X)$ 是闭且有界的. 我们从 $\mathbf{R}$ 的完备性推出实数的闭有界非空集有一个最大成员与一个最小成员.

为证明逆命题, 假定 $X$ 上的每个连续实值函数取到最大值与最小值. 根据定理 20, 为证明 $X$ 是紧的, 必要与充分的是证明它是全有界与完备的. 我们用反证法证明 $X$ 是全有界的. 若 $X$ 不是全有界的, 则存在 $r>0$ 与 $X$ 的可数无穷子集 $\{x_{n}\}_{n=1}^{\infty}$, 使得开球族 $\{B(x_{n},r)\}_{n=1}^{\infty}$ 是不交的. 对每个自然数 $n$, 定义函数 $f_{n}:X\to\mathbf{R}$ 为
\begin{align*}
f_{n}(x)=\begin{cases}
\frac{r}{2}-\rho(x,x_{n}) & \text{若 }\rho(x,x_{n})\leqslant\frac{r}{2}\\
0 & \text{否则}
\end{cases}
\end{align*}
定义函数 $f:X\to\mathbf{R}$ 为对所有 $x\in X$,
\begin{align*}
f(x)=\sum_{n=1}^{\infty}n\cdot f_{n}(x).
\end{align*}
由于每个 $f_{n}$ 连续, 在 $B(x_{n},r/2)$ 外等于零且 $\{B(x_{n},r)\}_{n=1}^{\infty}$ 是不交的, $f$ 是恰当定义与连续的. 但对每个自然数 $n$, $f(x_{n})=n\cdot r/2$, 因此 $f$ 无上界, 因此取不到最大值. 这是一个矛盾. 因此 $X$ 是全有界的. 剩下来要证明 $X$ 是完备的. 令 $\{x_{n}\}$ 为 $X$ 中的 Cauchy 序列. 接着对每个 $x\in X$, 我们从三角不等式推出 $\{\rho(x,x_{n})\}$ 是实数的 Cauchy 序列, 由于 $\mathbf{R}$ 是完备的, 该序列收敛到一个实数. 定义函数 $f:X\to\mathbf{R}$ 为
\begin{align*}
\text{对所有 }x\in X,\quad f(x)=\lim_{n\to\infty}\rho(x,x_{n}).
\end{align*}
又一次用三角不等式我们得出 $f$ 是连续的. 根据假设, 存在 $X$ 中的点 $x$ 使得 $f$ 在该点取到最小值. 由于 $\{x_{n}\}$ 是 Cauchy 的, $f$ 在 $X$ 上的最小值是 $0$. 于是 $f(x)=0$ 且因此 $\{x_{n}\}$ 收敛到 $x$. 因此 $X$ 是完备的.
\end{proof}

若 $\{\mathcal{O}_{\lambda}\}_{\lambda\in\Lambda}$ 是度量空间 $X$ 的开覆盖, 则每个点 $x\in X$ 包含于该覆盖的某个成员 $\mathcal{O}_{\lambda}$, 且由于 $\mathcal{O}_{\lambda}$ 是开的, 存在某个 $\varepsilon>0$, 使得
\begin{align}
B(x,\varepsilon)\subseteq\mathcal{O}_{\lambda}.\label{align:ch09-Lebesgue数}
\end{align}
一般来说, $\varepsilon$ 依赖于 $x$ 的选取. 下面的命题告诉我们, 对于紧度量空间这个包含关系在以下意义下一致成立: 能够找到与 $x\in X$ 无关的 $\varepsilon$ 使得\eqref{align:ch09-Lebesgue数}成立. 具有这个性质的正数 $\varepsilon$ 称为覆盖 $\{\mathcal{O}_{\lambda}\}_{\lambda\in\Lambda}$ 的 Lebesgue 数.

\begin{theorem}[Lebesgue 覆盖引理]\label{theorem:ch09-Lebesgue覆盖引理}
令 $\{\mathcal{O}_{\lambda}\}_{\lambda\in\Lambda}$ 为紧度量空间 $X$ 的开覆盖. 则存在数 $\varepsilon>0$, 使得对每个 $x\in X$, 开球 $B(x,\varepsilon)$ 包含于该覆盖的某个成员.
\end{theorem}

\begin{proof}
我们用反证法讨论. 假设没有这样的正 Lebesgue 数. 则对每个自然数 $n$, $1/n$ 不是一个 Lebesgue 数. 因此存在 $X$ 中的点, 记为 $x_{n}$, 使得 $B(x_{n},1/n)$ 不包含于该覆盖的某个成员. 这定义了 $X$ 中的一个序列 $\{x_{n}\}$. 根据紧性定理, $X$ 是序列紧的. 因此子序列 $\{x_{n_{k}}\}$ 收敛到点 $x_{0}\in X$. 现在存在某个 $\lambda_{0}\in\Lambda$ 使得 $\mathcal{O}_{\lambda_{0}}$ 包含 $x_{0}$, 且由于 $\mathcal{O}_{\lambda_{0}}$ 是开的, 存在中心在 $x_{0}$ 的球 $B(x_{0},r_{0})$, 使得
\begin{align*}
B(x_{0},r_{0})\subseteq\mathcal{O}_{\lambda_{0}}.
\end{align*}
我们可以选取指标 $k$ 使得 $\rho(x_{0},x_{n_{k}})<r_{0}/2$ 且 $1/n_{k}<r_{0}/2$. 根据三角不等式, $B(x_{n_{k}},1/n_{k})\subseteq\mathcal{O}_{\lambda_{0}}$, 而这与 $x_{n_{k}}$ 的选取使得 $B(x_{n_{k}},1/n_{k})$ 不包含于 $\{\mathcal{O}_{\lambda}\}_{\lambda\in\Lambda}$ 的某个成员矛盾.
\end{proof}

\begin{proposition}\label{proposition:ch09-紧集连续映射一致连续}
从紧度量空间 $(X,\rho)$ 到度量空间 $(Y,\sigma)$ 的连续映射是一致连续的.
\end{proposition}

\begin{proof}
令 $f$ 为从 $X$ 到 $Y$ 的连续映射. 令 $\varepsilon>0$. 根据在一个点连续的 $\varepsilon$-$\delta$ 准则, 对每个 $x\in X$, 存在 $\delta_{x}>0$ 使得若 $\rho(x,x')<\delta_{x}$, 则 $\sigma(f(x),f(x'))<\varepsilon/2$. 因此, 设 $\mathcal{O}_{x}=B(x,\delta_{x})$, 根据关于 $\sigma$ 的三角不等式,
\begin{align}
\sigma(f(u),f(v))\leqslant\sigma(f(u),f(x))+\sigma(f(x),f(v))<\varepsilon,\quad u,v\in\mathcal{O}_{x}.\label{align:ch09-紧一致连续估计}
\end{align}
令 $\delta$ 为开覆盖 $\{\mathcal{O}_{x}\}_{x\in X}$ 的 Lebesgue 数. 则对 $u$, $v\in X$, 若 $\rho(u,v)<\delta$, 存在某个 $x$ 使得 $u\in B(v,\delta)\subseteq\mathcal{O}_{x}$, 因此, 根据\eqref{align:ch09-紧一致连续估计}, $\sigma(f(u),f(v))<\varepsilon$.
\end{proof}

\begin{exercise}\label{exercise:ch09-52}
考虑由有理数组成的具有绝对度量的度量空间 $\mathbf{Q}$. $\mathbf{Q}$ 的什么子空间是完备的, 什么子空间是紧的?
\end{exercise}

\begin{exercise}\label{exercise:ch09-53}
令 $B=B(x,r)$ 为欧氏空间 $\mathbf{R}^{n}$ 的开球. 用以下步骤证明 $B$ 不是紧的: (i) 证明 $B$ 不是序列紧, (ii) 找到 $B$ 的一个没有有限子覆盖的开覆盖, (iii) 证明 $B$ 不是闭的.
\end{exercise}

\begin{exercise}\label{exercise:ch09-54}
什么时候具有离散度量的非空集 $X$ 是紧度量空间?
\end{exercise}

\begin{exercise}\label{exercise:ch09-55}
$\rho$ 和 $\sigma$ 是非空集 $X$ 上的等价度量. 证明度量空间 $(X,\rho)$ 是紧的当且仅当度量空间 $(X,\sigma)$ 是紧的.
\end{exercise}

\begin{exercise}\label{exercise:ch09-56}
证明两个紧度量空间的笛卡儿积也是紧的.
\end{exercise}

\begin{exercise}\label{exercise:ch09-57}
证明两个全有界的度量空间的笛卡儿积也是全有界的.
\end{exercise}

\begin{exercise}\label{exercise:ch09-58}
对于 $E$ 包含于度量空间 $X$, 证明子空间 $E$ 全有界当且仅当对每个 $\varepsilon>0$, $E$ 可被有限个中心在 $X$、半径为 $\varepsilon$ 的开球 (在 $X$ 中是开的) 覆盖.
\end{exercise}

\begin{exercise}\label{exercise:ch09-59}
令 $E$ 为紧度量空间 $X$ 的子集. 证明度量子空间 $E$ 是紧的当且仅当 $E$ 是 $X$ 的一个闭子集.
\end{exercise}

\begin{exercise}\label{exercise:ch09-60}
(Frechet 交定理) 令 $\{F_{n}\}_{n=1}^{\infty}$ 为紧度量空间 $X$ 的非空闭子集的下降可数族. 证明 $\bigcap_{n=1}^{\infty}F_{n}\neq\emptyset$.
\end{exercise}

\begin{exercise}\label{exercise:ch09-61}
对于度量空间 $X$ 的子集 $E$, 证明 $E$ 是全有界的当且仅当它的闭包 $\overline{E}$ 是全有界的.
\end{exercise}

\begin{exercise}\label{exercise:ch09-62}
对于完备度量空间 $X$ 的子集 $E$, 证明 $E$ 是全有界的当且仅当它的闭包 $\overline{E}$ 是紧的.
\end{exercise}

\begin{exercise}\label{exercise:ch09-63}
令 $B=\{\{x_{n}\}\in\ell^{2}\mid\sum_{n=1}^{\infty}x_{n}^{2}\leqslant 1\}$ 为 $\ell^{2}$ 的闭单位球. 用以下步骤证明 $B$ 不是紧的: (i) 证明 $B$ 不是序列紧, (ii) 找到 $B$ 的一个没有有限子覆盖的开覆盖, (iii) 证明 $B$ 不是全有界的.
\end{exercise}

\begin{exercise}\label{exercise:ch09-64}
令 $B=\{f\in L^{2}[a,b]\mid\|f\|_{2}\leqslant 1\}$ 为 $L^{2}[a,b]$ 的闭单位球. 用以下步骤证明 $B$ 不是紧的: (i) 证明 $B$ 不是序列紧, (ii) 找到 $B$ 的一个没有有限子覆盖的开覆盖, (iii) 证明 $B$ 不是全有界的.
\end{exercise}

\begin{exercise}\label{exercise:ch09-65}
令 $X$ 为全有界度量空间.
\begin{enumerate}[(i)]
\item 若 $f$ 是从 $X$ 到度量空间 $Y$ 的一致连续映射, 证明 $f(X)$ 是全有界的.
\item 若仅要求 $f$ 连续, (i) 仍然成立吗?
\end{enumerate}
\end{exercise}

\begin{exercise}\label{exercise:ch09-66}
令 $\rho$ 为集合 $X$ 上的度量. 定义
\begin{align*}
\tau(u,v)=\frac{\rho(u,v)}{1+\rho(u,v)},\quad\text{对所有 }u,v\in X.
\end{align*}
验证 $\tau$ 是 $X$ 上的有界度量且关于度量 $\rho$ 与关于度量 $\tau$ 的序列收敛性是相同的. 得出以下结论: 关于度量 $\rho$ 是闭的集合关于度量 $\tau$ 是闭的, 且关于度量 $\rho$ 是开的集合关于度量 $\tau$ 是开的. 度量 $\rho$ 和 $\tau$ 是等价的吗?
\end{exercise}

\begin{exercise}\label{exercise:ch09-67}
令 $E$ 为欧氏空间 $\mathbf{R}^{n}$ 的子集. 假设 $E$ 的每个连续实值函数取到最小值. 证明 $E$ 是闭且有界的.
\end{exercise}

\begin{exercise}\label{exercise:ch09-68}
令 $E$ 为欧氏空间 $\mathbf{R}^{n}$ 的子集. 假设 $E$ 的每个连续实值函数是一致连续的. 证明 $E$ 是闭且有界的.
\end{exercise}

\begin{exercise}\label{exercise:ch09-69}
假定 $f$ 是欧氏空间 $\mathbf{R}^{n}$ 上的连续实值函数, 其具有以下性质: 存在一个数 $c$ 使得对所有 $x\in\mathbf{R}^{n}$, $|f(x)|\geqslant c\cdot\|x\|$. 证明: 若 $K$ 是实数的紧集, 则它在 $f$ 的原象 $f^{-1}(K)$ 也是紧的. (具有该性质的映射称为适当的.)
\end{exercise}

\begin{exercise}\label{exercise:ch09-70}
对紧度量空间 $(X,\rho)$, 证明存在点 $u$, $v\in X$ 使得 $\rho(u,v)=\mathrm{diam}X$.
\end{exercise}

\begin{exercise}\label{exercise:ch09-71}
令 $K$ 为度量空间 $(X,\rho)$ 的紧子集而 $x_{0}$ 属于 $X$. 证明存在点 $z\in K$ 使得对所有 $x\in K$, $\rho(z,x_{0})\leqslant\rho(x,x_{0})$.
\end{exercise}

\begin{exercise}\label{exercise:ch09-72}
令 $K$ 为度量空间 $X$ 的紧子集. 对于点 $x\in X\sim K$, 证明存在包含 $K$ 的开集 $\mathcal{U}$ 和包含 $x$ 的开集 $\mathcal{O}$ 使得 $\mathcal{U}\bigcap\mathcal{O}=\emptyset$.
\end{exercise}

\begin{exercise}\label{exercise:ch09-73}
令 $A$ 和 $B$ 为度量空间 $(X,\rho)$ 的子集. 定义
\begin{align*}
\mathrm{dist}(A,B)=\inf\{\rho(u,v)\mid u\in A,v\in B\}.
\end{align*}
若 $A$ 是紧的而 $B$ 是闭的, 证明 $A\bigcap B=\emptyset$ 当且仅当 $\mathrm{dist}(A,B)>0$.
\end{exercise}

\begin{exercise}\label{exercise:ch09-74}
令 $K$ 为度量空间 $X$ 的紧子集而 $\mathcal{O}$ 是一个包含 $K$ 的开集. 利用前一个习题证明存在一个开集 $\mathcal{U}$ 满足
\begin{align*}
K\subseteq\mathcal{U}\subseteq\overline{\mathcal{U}}\subseteq\mathcal{O}.
\end{align*}
\end{exercise}

\end{document}